 \setcounter{subsection}{0}

\section{CONCLUSIONES}

\subsection{Análisis de Impacto del Proyecto}

\subsubsection{Impacto en la ingeniería de sistemas}

Este capítulo explica qué dejó el trabajo a la ingeniería de sistemas y a su contexto, comprueba el cumplimiento de cada objetivo específico, delimita el alcance de la evidencia y propone el trabajo futuro. Además del sistema construido, el trabajo deja tres prácticas que se pueden trasladar a otros proyectos.

Primero, la arquitectura se eligió comparando alternativas con el método ATAM \citep{bass_software_2021}, apoyado en impulsores comprobables: recuperabilidad asimétrica, frontera de cumplimiento y perfil de recurso heterogéneo.

Segundo, el sistema se auditó más allá de sus requisitos funcionales, con revisiones de autorización (22 de 152 puntos de entrada con comprobación incompleta) y de datos personales (catorce hallazgos), cuyos pendientes reúne la Sección~\ref{sec:defectos}.

Tercero, las reglas sobre qué servicio puede usar a cuál se volvieron pruebas automáticas llamadas funciones de aptitud \citep{ford_building_2017}: doce comparan el código con el manifiesto (el archivo que declara el lenguaje, la base y las tablas de cada servicio) y otras veintiocho comprueban que detectan infracciones plantadas.

\subsubsection{Impacto en contextos globales, económicos, ambientales y sociales}

Según la encuesta del OE1 (Capítulo~IV), el 65\,\% de las familias dedica cuatro horas o más a comprar los útiles de un hijo. El OE5 (Sección~\ref{sec:oe5}) midió si Tukit lo reduce: la carga percibida fue de 9, frente a 63 recordados para su última compra (RTLX, de 0 a 100), y la usabilidad (SUS, de 0 a 100) promedió 90, con seis personas. Solo uno de cinco eligió su kit sin ayuda.

En lo económico, el modelo de negocio adoptado \citep{osterwalder_business_2010} es B2B2C (de empresa a colegio y de ahí al consumidor): el colegio no paga y el ingreso viene de la compra del acudiente. Su viabilidad financiera se analiza en el Capítulo~VII, con los supuestos que delimita la Sección~\ref{sec:amenazas}.

En el plano global, los resultados solo valen para Colombia, porque el calendario escolar y la normativa de las listas son locales y el trabajo no probó otro país. En el plano tecnológico, una lista escolar en foto o PDF se convierte en un pedido pagado con un modelo multimodal, un catálogo curado y un optimizador. En el plano ambiental no se midió un efecto directo; que el kit reduzca viajes y compras redundantes es solo una hipótesis. En lo social, la encuesta se aplicó en un solo colegio del sector Salitre y no se generaliza a la población escolar colombiana.

\subsection{Conclusiones y Trabajo Futuro}

Los cinco objetivos específicos se cumplieron, cada uno con el método de su fase. La Tabla~\ref{tab:cumplimiento} resume, por objetivo, la evidencia, su sección, el estado y el alcance; los párrafos siguientes dan la razón de cada uno.

\begin{table}[!htb]
\centering\footnotesize
\setlength{\tabcolsep}{3pt}
\caption{Cumplimiento por objetivo específico.}
\label{tab:cumplimiento}
\begin{tabularx}{\textwidth}{@{}p{0.9cm}Xp{2.9cm}p{1.9cm}p{3.2cm}@{}}
\toprule
OE & Evidencia & Sección & Estado & Alcance \\
\midrule
OE1 & Corpus de listas, conceptos canónicos y encuesta (n=60) & Cap.~III y IV & Cumplido & Encuesta en un solo colegio \\
OE2 & Comparación ATAM de tres alternativas, vistas 4+1 y registro de decisiones & Cap.~V & Cumplido & Puntajes de la matriz ATAM asignados por el equipo \\
OE3 & F1 micro de 0,948 en lectura; resolución con $\tau=0{,}75$ (precisión \calTauP{}); NL2C con 58 de 70 tareas y 0 de 81 ataques exitosos & Sec.~\ref{sec:percepcion}, \ref{sec:optimizacion-cal} y~\ref{sec:nl2c} & Cumplido & Un revisor; oro corregido sobre Gemini; ventaja amplia sobre las reglas y menor sobre Vision \\
OE4 & Optimizador de cuatro objetivos, MAUT con pesos de la P21, hipervolumen de 0,9969 del exacto y aplicación web con pago en modo simulado & Sec.~\ref{sec:optimizacion} y~\ref{sec:optimizacion-cal} & Cumplido & Pago simulado; una semilla \\
OE5 & Pruebas de unidad, integración, sistema y carga; ESC-D4 con p95 de 0,31~s frente a 8~s; seis sesiones con usuarios reales, los seis con pedido pagado; criterios C1 y C2 cumplidos y C3 en el pago (T6) & Sec.~\ref{sec:control-calidad}, \ref{sec:optimizacion} y~\ref{sec:oe5} & Cumplido & Seis acudientes; C3 formativo en T1 y T4 \\
\bottomrule
\end{tabularx}
\end{table}

El OE1 (fase F1, Design Thinking en sus etapas de empatizar y definir, con la norma IEEE 830 para el SRS) se cumplió porque la encuesta, las entrevistas y la caracterización del corpus permitieron definir un conjunto de conceptos canónicos lo bastante reducido para que un catálogo finito lo represente (Capítulos~III y~IV).

Con ese corpus como entrada, el OE2 (fase F2, con el proceso RUP, la norma IEEE 1016 para el documento de diseño de software, SDD, y ATAM) se cumplió porque entregó un diseño del sistema documentado y justificado, descrito con vistas 4+1 \citep{kruchten_41_1995} y la norma ISO/IEC/IEEE 42010 \citep{iso_42010_2022}.

Sobre ese diseño, el OE3 (SCRUM por sprints y CRISP-DM en sus fases de comprensión, preparación, modelado y evaluación) cubrió dos tareas: la percepción, que es leer los renglones de una lista, y la resolución, que es emparejarlos con un producto del catálogo. En percepción, el modelo multimodal superó al reconocimiento con reglas (F1 micro de 0,948 frente a 0,298). En resolución, una sola fase suma el parecido léxico y el coseno con $\tau=0{,}75$, y se compensa con la revisión de cada lista nueva antes de publicarla.

Con los renglones resueltos, el OE4 (SCRUM y CRISP-DM aplicados a la optimización) se cumplió con un optimizador de cuatro objetivos (costo, fidelidad de marca, disponibilidad y personaje). Hasta 4.096 combinaciones las enumera todas y obtiene el frente exacto (los kits que no pueden mejorar en un objetivo sin empeorar otro); por encima recurre a NSGA-II \citep{deb_fast_2002,blank_pymoo_2020}, con un hipervolumen de 0,9969 del exacto (Sección~\ref{sec:optimizacion-cal}). El personaje es un hallazgo del desarrollo: la encuesta midió esa preferencia y los tres criterios de la VFP no la recogían. El prototipo cubre el recorrido de la lista al pedido pagado.

Para medir ese recorrido, el OE5 (fase F5, con un protocolo preregistrado con ISO/IEC 25010) pedía evaluar el sistema con pruebas funcionales, no funcionales, de carga y de usabilidad con usuarios reales, y se cumplió con el control de calidad en cuatro niveles, la prueba de carga, la prueba ESC-D4 y seis sesiones con acudientes del Colegio Agustiniano Ciudad Salitre. Los criterios preregistrados C1 y C2 (SUS y carga percibida) se cumplieron, y C3 (éxito sin ayuda) se alcanzó en el pago (T6); el registro (T1) y la elección del kit (T4) se leen como formativos (Sección~\ref{sec:oe5}). Cumplido significa que se ejecutaron las cuatro clases de prueba y se midió la experiencia, no que el sistema esté listo para operar; por eso el trabajo futuro empieza por producción.

El objetivo general, diseñar el sistema de recomendación multiobjetivo y medir su efecto con padres del sector Salitre, se da por cumplido porque los OE1 a OE4 produjeron el sistema y el OE5 lo probó con acudientes reales. Los seis llegaron a un pedido pagado (simulado), los criterios preregistrados C1 (SUS de 90, sobre el umbral de 68) y C2 (diferencia de carga percibida que excluye el cero, con $d_z \ge 0{,}5$) se cumplieron y la satisfacción se describe con el NPS de $+83$; con seis participantes los efectos son formativos, dentro del alcance que fija la subsección siguiente.

\subsubsection{Limitaciones}
\label{sec:amenazas}

Las conclusiones anteriores valen dentro del alcance que sigue, ordenado con los tipos de validez de Wohlin et al. \citep{wohlin_experimentation_2012}. Las corridas de $\tau$ y la evaluación con usuarios usaron la misma versión del código, lo que permite reproducirlas.

\textbf{Medición.} Solo una persona anotó el corpus, así que el kappa de Cohen \citep{cohen_coefficient_1960} compara una heurística con un modelo de lenguaje y no a dos anotadores humanos ciegos; el 0,7331 es entre dos métodos automáticos. El oro de percepción se construyó corrigiendo la salida de Gemini, con la que coincide en el 86,7\,\% de los ítems literales, de modo que el 0,948 favorece a ese motor. Un solo revisor etiquetó los \calTauPares{} pares de $\tau$, \calTauColegioMayor{} de ellos de un mismo colegio, y su puntaje toma solo \calTauValores{} valores distintos. Los pesos del optimizador salen de lo que los acudientes dicen preferir y no de lo que compran.

\textbf{Percepción y latencias.} La percepción se corrió una vez por motor sobre diez documentos y sin versión fija del modelo remoto. Su latencia con el modelo real (ESC-D1) y la de la recuperación de candidatos en el sistema construido no se midieron, porque en las pruebas de carga la percepción fue simulada. Las cifras del banco de alternativas y de la programación dinámica son de una sola semilla, y el peor caso de la caída del catálogo se acotó con un modelo (300,8 s) y no se midió.

\textbf{Pruebas del sistema.} Los niveles 2 y 3 de caos no se repitieron completos: la campaña del 5 de noviembre ejecutó 25 ensayos y omitió 14 por decisión del equipo. Las fallas CH-09, CH-24 y CH-30 quedan sin arreglo, y las correcciones de octubre están en una rama aparte, sin integrar. La telemetría no cumple su umbral. En WSL2, tres escenarios de carga (L-03, L-05 y L-06) no cumplen y las cargas de Windows y Mac no son comparables; la carga y los tiempos del optimizador son de un solo equipo. El pago se probó con la pasarela simulada y no con las llaves de pruebas de Wompi, y no se probó en dispositivos reales; tampoco se ejecutaron la tabla de niveles de usabilidad ni la auditoría de dependencias de Python. En el asistente, el criterio para descartar casos de la última ronda se cambió sin preregistro; en el OE5, la tarea T8 se agregó después del preregistro y es exploratoria, y ningún acudiente usó el asistente, así que su utilidad no está medida.

\textbf{Muestra y estadística.} La encuesta, 391 de los 405 pares de $\tau$, las auditorías de los kits y las seis sesiones vienen del sector Salitre, con el estrato 4 sobrerrepresentado (68,3\,\%) y muestreo por conveniencia; la prueba de la calibración tuvo solo \calColegiosPrueba{} colegios. Con menos de diez participantes la potencia no alcanza para detectar un efecto de tamaño medio, por lo que el SUS, el RTLX y el éxito por tarea se leen como formativos: C1 y C2 se cumplen en esta muestra, pero no se extienden a la población de acudientes.

\textbf{Proyección financiera.} Descansa en supuestos del equipo (volúmenes, costos fijos, inversión, comisión, logística y conversión) y el ticket sale de un solo colegio; con una comisión del 10\,\% la contribución por kit es negativa. Son cifras que la operación real debe confirmar.

\textbf{Alcance del código y del asistente.} El nombre del menor, si figura en el documento cargado, viaja en la imagen hacia Gemini y queda 30 días en la lectura guardada en Redis (PE-06). Los índices HNSW y de trigramas existen pero no están activos, y la búsqueda es en memoria hasta cerca de 50.000 productos (ESC-E4). El corte del 10\,\% de sobreprecio para marcar el Premium se fijó sin medirlo. Del asistente en lenguaje natural solo está validada su seguridad: ningún acudiente lo usó en el OE5, así que su utilidad no se midió.

\subsubsection{Trabajo futuro}
\label{sec:trabajo-futuro}

El trabajo futuro sigue las cinco líneas de la propuesta de grado (VFP~v4), más el lanzamiento comercial.
\begin{enumerate}
\item \textbf{Preparar el sistema para producción.} Es la condición para operar con familias reales: servicio del catálogo con varios procesos, que reduce la latencia bajo carga, y cierre de los 85 hallazgos vigentes del control de calidad, 16 de severidad alta (Sección~\ref{sec:calidad-resumen}): autorización uniforme en las rutas, rol de acceso por base de datos, depuración del nombre del menor en el documento cargado, supresión de datos, medición de la cobertura en la integración continua y cobro real con la pasarela.
\item \textbf{Cambios de la prueba con usuarios.} Son los ajustes a la tarea de elección del kit (T4 de las sesiones del OE5), el panel de colegio (requerimientos CO-01 y CO-02) y el ajuste de presupuesto en vivo.
\item \textbf{Alianza con tiendas de uniformes e integración de libros.} Acordar con las tres tiendas de uniformes y conectar Buscalibre y las editoriales; el plan lector está poco cubierto (18 de 78 líneas de libros tienen su edición en el catálogo).
\item \textbf{Calibración con más anotadores y más colegios.} Repetir la calibración del umbral y de los pesos con dos anotadores y más colegios.
\item \textbf{Lanzamiento comercial.} Lanzar en el colegio piloto, con la institución como canal de entrada, en la temporada de diciembre a enero y con referidos entre padres. Cuatro indicadores (conversión de lista cargada a kit comprado, costo de adquisición, activación de los colegios y recompra) reemplazan los supuestos del modelo financiero por mediciones reales.
\item \textbf{Alcance del producto.} Aplicaciones nativas con modo sin conexión, rutas de entrega agrupadas entre tiendas y reventa de textos entre hogares.
\end{enumerate}
